\documentclass[11pt,letterpaper]{article}
\usepackage[letterpaper,left=0.87in,right=0.87in,top=0.72in,bottom=0.72in,headheight=14pt,headsep=15pt,footskip=25pt]{geometry}
\usepackage[T1]{fontenc}
\usepackage[utf8]{inputenc}
\usepackage{newtxtext,newtxmath}
\usepackage{microtype,tabularx,array,enumitem,titlesec,needspace,fancyhdr}
\usepackage[hidelinks]{hyperref}
\hypersetup{pdfauthor={Dingli Wang},pdfcreator={LaTeX},pdflang={en-US}}
\urlstyle{same}
\setlength{\parindent}{0pt}
\setlength{\parskip}{6.0pt}
\setlength{\emergencystretch}{1.5em}
\widowpenalty=10000
\clubpenalty=10000
\displaywidowpenalty=10000
\raggedbottom
\pagestyle{fancy}
\fancyhf{}
\renewcommand{\headrulewidth}{0pt}
\renewcommand{\footrulewidth}{0pt}
\lhead{\small Dingli Wang}
\rhead{\small \DocName}
\cfoot{\small\thepage}
\titleformat{\section}{\normalsize\bfseries}{}{0pt}{}
\titlespacing*{\section}{0pt}{10pt}{4pt}
\newcommand{\docheader}[1]{%
  \thispagestyle{plain}%
  {\Large\bfseries Dingli Wang}\hfill{\large #1}\par
  \vspace{2pt}
  {\small Department of Economics, University of California, Riverside}\par
  {\small\href{mailto:dingli.wang@email.ucr.edu}{dingli.wang@email.ucr.edu}\quad\textbar\quad
  \href{https://dingliwang.github.io/}{dingliwang.github.io}\quad\textbar\quad\href{https://papers.ssrn.com/sol3/cf_dev/AbsByAuth.cfm?per_id=10994248}{SSRN}}\par
  \vspace{4pt}\hrule height 0.4pt\vspace{9pt}
}
\newcommand{\coauthor}[1]{{\normalfont (with #1)}}

\newcommand{\paperlinks}[2]{{\fontsize{10.5}{12}\selectfont\normalfont\href{#1}{\underline{Paper}}\hspace{1.0em}\textbar\hspace{1.0em}\href{#2}{\underline{Slides}}}}

\newcommand{\DocName}{Curriculum Vitae}
\hypersetup{pdftitle={Dingli Wang - Curriculum Vitae}}
\setlength{\parskip}{1pt}
\renewcommand{\baselinestretch}{1.0}
\titleformat{\section}{\normalsize\scshape}{}{0pt}{}[{\vspace{1pt}\titlerule[0.35pt]}]
\titlespacing*{\section}{0pt}{7pt}{4pt}
\newcommand{\cvrow}[2]{\begin{tabularx}{\textwidth}{@{}Xr@{}}#1&#2\end{tabularx}\par\vspace{0.5pt}}
\newcommand{\paper}[3]{\needspace{4\baselineskip}\textbf{#1}\par {#2}\par {\small #3}\par}
\lfoot{\small Updated November 2026}
\cfoot{}
\rfoot{\small\thepage}
\fancypagestyle{plain}{\fancyhf{}\lfoot{\small Updated November 2026}\rfoot{\small\thepage}\renewcommand{\headrulewidth}{0pt}}
\begin{document}
\thispagestyle{plain}
{\fontsize{25}{28}\selectfont\bfseries Dingli Wang}\par
\vspace{1pt}
Department of Economics, University of California, Riverside\par
\href{mailto:dingli.wang@email.ucr.edu}{dingli.wang@email.ucr.edu}\quad\textbar\quad (951)~231--0895\quad\textbar\quad
\href{https://dingliwang.github.io/}{dingliwang.github.io}\quad\textbar\quad\href{https://papers.ssrn.com/sol3/cf_dev/AbsByAuth.cfm?per_id=10994248}{SSRN}

\vspace{1pt}
{\small
\begin{tabularx}{\textwidth}{@{}X X@{}}
\textbf{Placement Director} & \textbf{Placement Coordinator}\\
Joseph Cummins & Gary Kuzas\\
(951)~827--1582 \quad \href{mailto:joseph.cummins@ucr.edu}{joseph.cummins@ucr.edu} &
(951)~827--1474 \quad \href{mailto:gary.kuzas@ucr.edu}{gary.kuzas@ucr.edu}
\end{tabularx}}

\section{Education}
\cvrow{\textbf{University of California, Riverside} --- Ph.D. in Economics}{Expected June 2027}
{\small Dissertation committee: Tae-Hwy Lee (Chair), Donggyu Kim, Ruoyao Shi.}\par
\vspace{1pt}
\cvrow{\textbf{Boston University} --- M.A. in Economics}{2020}
\vspace{1pt}
\cvrow{\textbf{Tianjin University of Finance and Economics} --- B.S. in Economic Statistics}{2018}

\section{Research Fields}
Econometrics; Financial Econometrics; Forecasting; Machine Learning and AI in Economics

\section{Job Market Papers}

\needspace{6\baselineskip}
\href{https://dingliwang.github.io/papers/Robust-Replacement-of-Human-Decision-Makers-under-Distribution-Shift.pdf}{\textbf{Robust Replacement of Human Decision Makers under Distribution Shift}}\par
\begin{tabularx}{\textwidth}{@{}Xr@{}}
  \textit{Revise and Resubmit, Journal of Econometrics} &
  \paperlinks{https://dingliwang.github.io/papers/Robust-Replacement-of-Human-Decision-Makers-under-Distribution-Shift.pdf}{https://dingliwang.github.io/papers/Robust-Replacement-of-Human-Decision-Makers-under-Distribution-Shift-Slides.pdf}
\end{tabularx}
{\fontsize{11}{15}\selectfont This paper studies when a machine learning algorithm should replace a human decision maker under distribution shift. We infer each human's error rates from past decisions, assuming that they stay the same on future cases, while the data seen by the algorithm may shift within a Wasserstein ball. For each threshold, we compute the algorithm's worst false positive and true positive rates over the ball, which gives a worst-case ROC curve. We replace a human only if some point on this curve is, with high posterior probability, at least as good as the human in both rates. We show that the estimated curve converges uniformly and bound the welfare loss from replacement. We also compare the rule with an infeasible benchmark that knows the deployment distribution and each human's error rates and sets the algorithm's false positive rate equal to the human's. Relative to this benchmark, choosing the threshold with the highest posterior probability can leave a welfare loss that does not vanish with more data. In simulations, robust rules replace fewer humans than a rule that ignores shifts, giving up some gains without shifts but losing less after large shifts.\par}

\vspace{2pt}
\needspace{6\baselineskip}
\href{https://dingliwang.github.io/papers/Median-Anchored-Adjustment-of-Joint-VaR-ES-Forecasts.pdf}{\textbf{Median-Anchored Adjustment of Joint VaR--ES Forecasts}} {\small\normalfont (with Tae-Hwy Lee)}\par
\begin{tabularx}{\textwidth}{@{}Xr@{}}
  \textit{Under Review, Journal of Business \& Economic Statistics} &
  \paperlinks{https://dingliwang.github.io/papers/Median-Anchored-Adjustment-of-Joint-VaR-ES-Forecasts.pdf}{https://dingliwang.github.io/papers/Median-Anchored-Adjustment-of-Joint-VaR-ES-Forecasts-Slides.pdf}
\end{tabularx}
{\fontsize{11}{15}\selectfont Risk managers often work with a fitted forecasting model they cannot replace even when its tail forecasts need adjustment. We propose a median-anchored rule that multiplies the distances from the fitted median to Value-at-Risk (VaR) and Expected Shortfall (ES) by a common positive multiplier. The rule preserves VaR--ES ordering, and minimizing VaR check loss gives a closed-form weighted-quantile estimator. We establish consistency, give an asymptotic distribution under high-level conditions accounting for baseline estimation, and derive a VaR coverage-error bound at the forecast origin. When the same multiplier correctly adjusts both VaR and ES, the adjusted pair has lower conditional zero-homogeneous Fissler--Ziegel (FZ0) risk at that origin. Monte Carlo experiments examine this condition and several forms of misspecification. In rolling S\&P~500 forecasts the adjustment lowers FZ0 loss in all four GARCH cells, each pairwise significant, with the largest and most robust gain, 10.4\%, for Gaussian GARCH at the 1\% tail; the 5\% gains do not survive the most conservative family-wise adjustments. A quantile-regression baseline marks the boundary: when the fitted tail is already flexible, one multiplier adds little. Filtered historical simulation quantifies what standardized residuals and conditional scales add when available.\par}

\newpage
\section{Other Research Papers}

\needspace{3\baselineskip}
\textbf{When Should Forecasts Be Equally Weighted? Signal Recovery and Estimation Costs}\par
\textit{Under Review, International Journal of Forecasting}\par
{\small Compares equal and estimated forecast-combination weights through signal recovery and estimation costs.}\par

\vspace{1pt}
\needspace{3\baselineskip}
\textbf{Adversarial Density Forecast for Macro-financial Risks} {\small\normalfont (with Tae-Hwy Lee)}\par
{\small Develops conditional generative density forecasts for short macroeconomic samples.}\par

\section{Conference Presentations and Seminars}
\cvrow{Econometrics Seminar, UC Riverside}{November 2026}
\cvrow{Midwest Econometrics Group (MEG), Cincinnati}{October 2026}
\cvrow{International Association for Applied Econometrics (IAAE), Lisbon}{June 2026}
\cvrow{Asian Meeting of the Econometric Society (AMES), Hong Kong}{June 2026}
\cvrow{Econometrics Seminar, UC Riverside}{April 2026}

\section{Teaching Experience}
\textit{University of California, Riverside}\par
\vspace{2pt}
\textbf{Instructor of Record}\par
\cvrow{ECON 182 --- Trade, Globalization, and Development}{Summer 2026}
\cvrow{ECON/BUS 178 --- International Trade}{Summer 2025}
\cvrow{\textbf{Instructor} --- Statistics Camp (Ph.D.)}{Summer 2026}
\vspace{1pt}
\textbf{Teaching Assistant}\par
\cvrow{ECON 205A --- Econometric Methods I (Ph.D.)}{Fall 2026}
\cvrow{ECON 205C --- Econometric Methods III (Ph.D.)}{Spring 2026}
\cvrow{ECON 107 --- Introductory Econometrics I}{Spring 2025}
\cvrow{ECON 101 --- Statistics for Economics}{2023--2026 (8 terms)}
\cvrow{ECON 003 --- Introduction to Microeconomics}{Winter 2023}
\cvrow{ECON 002 --- Introduction to Macroeconomics}{2022--2024 (4 terms)}

\section{Fellowships and Awards}
\cvrow{Outstanding Teaching Assistant Award, University of California, Riverside}{2026}
\cvrow{Dean's Distinguished Fellowship, University of California, Riverside}{2021--2026}
\cvrow{High Pass, Ph.D. Qualifying Examination, University of California, Riverside}{2022}
\cvrow{First Prize, Contemporary Undergraduate Mathematical Contest in Modeling}{2016}

\section{Computing and Languages}
\textbf{Computing:} Python, R, Stata, EViews, \LaTeX\par
\textbf{Languages:} English (fluent); Mandarin Chinese (native)


\vspace{-3pt}
\section{References\hfill{\normalfont\footnotesize University of California, Riverside}}

\begingroup
\setlength{\tabcolsep}{0pt}
\renewcommand{\arraystretch}{1}
\begin{tabularx}{\textwidth}{@{}
>{\raggedright\arraybackslash}X
@{\hspace{1.3em}}
>{\raggedright\arraybackslash}X@{}}
\begin{minipage}[t]{\linewidth}\vspace{0pt}\footnotesize\setlength{\parskip}{0pt}
{\small\textbf{\href{http://faculty.ucr.edu/~taelee/}{Tae-Hwy Lee}} (Chair)}\par
Professor\par
\href{mailto:tae.lee@ucr.edu}{tae.lee@ucr.edu}\par
(951) 827-1509
\end{minipage}
&
\begin{minipage}[t]{\linewidth}\vspace{0pt}\footnotesize\setlength{\parskip}{0pt}
{\small\textbf{\href{https://sites.google.com/site/donggyukim0329/home}{Donggyu Kim}}}\par
Professor\par
\href{mailto:donggyu.kim@ucr.edu}{donggyu.kim@ucr.edu}\par
(951) 827-8372
\end{minipage} \\[3pt]
\begin{minipage}[t]{\linewidth}\vspace{0pt}\footnotesize\setlength{\parskip}{0pt}
{\small\textbf{\href{https://ruoyaoshi.github.io}{Ruoyao Shi}}}\par
Associate Professor\par
\href{mailto:ruoyao.shi@ucr.edu}{ruoyao.shi@ucr.edu}\par
(951) 827-1494
\end{minipage}
&
\begin{minipage}[t]{\linewidth}\vspace{0pt}\footnotesize\setlength{\parskip}{0pt}
{\small\textbf{\href{https://sites.google.com/site/vsovero/home}{Veronica T. Sovero}}}\par
Associate Professor of Teaching\par
\href{mailto:veronica.sovero@ucr.edu}{veronica.sovero@ucr.edu}\par
(951) 827-0224
\end{minipage}
\end{tabularx}
\endgroup

\end{document}
